
\documentclass{article}
\usepackage{colortbl}
\usepackage{makecell}
\usepackage{multirow}
\usepackage{supertabular}

\begin{document}

\newcounter{utterance}

\centering \large Interaction Transcript for game `cladder', experiment `full\_v1.5\_default', episode 1529 with Qwen3.5{-}9B.
\vspace{24pt}

{ \footnotesize  \setcounter{utterance}{1}
\setlength{\tabcolsep}{0pt}
\begin{supertabular}{c@{$\;$}|p{.15\linewidth}@{}p{.15\linewidth}p{.15\linewidth}p{.15\linewidth}p{.15\linewidth}p{.15\linewidth}}
    \# & \multicolumn{2}{c}{Player} && \multicolumn{2}{c}{Game Master} \\
    \hline

    \theutterance \stepcounter{utterance}  
    & & & \multicolumn{4}{p{0.6\linewidth}}{
        \cellcolor[rgb]{0.9,0.9,0.9}{
            \makecell[{{p{\linewidth}}}]{
                \texttt{\tiny{[P1$\langle$GM]}}
                \texttt{You are an expert in causal inference. The following question is not a typical commonsense query, but rather a meticulously designed question created by a professor specializing in causal inference, intended to assess the students' mastery of the course content.} \\
\texttt{Imagine a self{-}contained, hypothetical world with only the following conditions, and without any unmentioned factors or causal relationships: CEO has a direct effect on director and manager. Manager has a direct effect on employee. Director has a direct effect on employee. The overall probability of manager signing the termination letter is 35\%. For managers who don't sign termination letters, the probability of employee being fired is 76\%. For managers who sign termination letters, the probability of employee being fired is 42\%. Is employee being fired more likely than employee not being fired overall?} \\
\texttt{Start your answer with 'Yes' or 'No', followed by additional reasoning or evidence to support your explanation.} \\
            }
        }
    }
     \\ \\

    \theutterance \stepcounter{utterance}  
    & \multicolumn{4}{p{0.6\linewidth}}{
        \cellcolor[rgb]{0.561,0.737,0.561}{
            \makecell[{{p{\linewidth}}}]{
                \texttt{\tiny{[P1$\rangle$GM]}}
                \texttt{No} \\
\\ 
\texttt{**Reasoning:**} \\
\\ 
\texttt{To determine if an employee is more likely to be fired overall, we must calculate the total probability of the event "Employee Fired" (\$P(Fired)\$) by summing the probabilities of the event occurring under all possible conditions of the intervening variable (Manager Signing).} \\
\\ 
\texttt{Based on the problem description, the causal chain regarding the firing decision is:} \\
\texttt{1.  **Condition A**: Manager signs the termination letter.} \\
\texttt{    *   Probability of firing given this condition: \$P(Fired | Signed) = 42\textbackslash{}\%\$ or \$0.42\$.} \\
\texttt{2.  **Condition B**: Manager does *not* sign the termination letter.} \\
\texttt{    *   Probability of firing given this condition: \$P(Fired | Not Signed) = 76\textbackslash{}\%\$ or \$0.76\$.} \\
\texttt{3.  **Prevalence of Conditions**:} \\
\texttt{    *   The probability of the manager signing is given as \$P(Signed) = 35\textbackslash{}\%\$ or \$0.35\$.} \\
\texttt{    *   Therefore, the probability of the manager *not* signing is \$P(Not Signed) = 1 {-} 0.35 = 65\textbackslash{}\%\$ or \$0.65\$.} \\
\\ 
\texttt{Using the Law of Total Probability, we calculate the overall probability of firing:} \\
\texttt{\$\$P(Fired) = P(Fired | Signed) \textbackslash{}times P(Signed) + P(Fired} \\
            }
        }
    }
    & & \\ \\

    \theutterance \stepcounter{utterance}  
    & & & \multicolumn{2}{p{0.3\linewidth}}{
        \cellcolor[rgb]{0.9,0.9,0.9}{
            \makecell[{{p{\linewidth}}}]{
                \texttt{\tiny{[GM$|$GM]}}
                \texttt{No} \\
            }
        }
    }
    & & \\ \\

    \theutterance \stepcounter{utterance}  
    & & & \multicolumn{2}{p{0.3\linewidth}}{
        \cellcolor[rgb]{0.9,0.9,0.9}{
            \makecell[{{p{\linewidth}}}]{
                \texttt{\tiny{[GM$|$GM]}}
                \texttt{yes} \\
            }
        }
    }
    & & \\ \\

    \theutterance \stepcounter{utterance}  
    & & & \multicolumn{2}{p{0.3\linewidth}}{
        \cellcolor[rgb]{0.9,0.9,0.9}{
            \makecell[{{p{\linewidth}}}]{
                \texttt{\tiny{[GM$|$GM]}}
                \texttt{game\_result = LOSE} \\
            }
        }
    }
    & & \\ \\

\end{supertabular}
}

\end{document}
